\section{Aplikasi dan Studi Kasus}

Bagian ini menyajikan dua studi kasus yang mengaitkan proposisi dan operator logika dengan spesifikasi sistem serta aturan bisnis, dilengkapi contoh implementasi singkat dalam Python \cite{levin_discrete,python_docs_boolean,rosen2019}.

\subsection*{Studi Kasus 1: Sistem Lampu Peringatan}

Sebuah sistem lampu peringatan di laboratorium diaktifkan berdasarkan kondisi berikut: \emph{lampu menyala} jika dan hanya jika \textbf{(sensor asap mendeteksi asap) ATAU (sensor suhu mendeteksi suhu tinggi)}. Kita dapat memodelkan dengan proposisi:
\begin{itemize}
  \item $p$: Sensor asap mendeteksi asap
  \item $q$: Sensor suhu mendeteksi suhu tinggi
  \item $L$: Lampu peringatan menyala
\end{itemize}
Spesifikasi: $L \leftrightarrow (p \vee q)$. Artinya, lampu menyala tepat ketika minimal salah satu dari $p$ atau $q$ benar. Tabel kebenaran untuk $p \vee q$ (yang ekuivalen dengan kondisi ``lampu harus menyala'') memiliki 4 baris: $(p,q)$ = (B,B), (B,S), (S,B), (S,S); hasil $p \vee q$ berturut-turut B, B, B, S. Jadi lampu padam hanya ketika kedua sensor tidak mendeteksi bahaya.

\begin{contoh}
\textbf{Implementasi kondisi lampu dalam Python}

\begin{lstlisting}[style=pythonstyle, caption={Sistem lampu peringatan: kondisi Boolean}]
def lampu_menyala(sensor_asap, sensor_suhu):
    """Lampu menyala jika asap TERDETEKSI atau suhu tinggi."""
    return sensor_asap or sensor_suhu

# Simulasi
print(lampu_menyala(True, False))   # True  -> lampu menyala
print(lampu_menyala(False, True))  # True  -> lampu menyala
print(lampu_menyala(False, False)) # False -> lampu padam
\end{lstlisting}
\end{contoh}

\subsection*{Studi Kasus 2: Kondisi Kelulusan Mahasiswa}

Aturan akademik: \emph{Mahasiswa dinyatakan lulus} jika dan hanya jika \textbf{IPK $\geq 2{,}0$ DAN tidak ada nilai D}. Misalkan $p$ = ``IPK $\geq 2{,}0$'', $q$ = ``Tidak ada nilai D''. Maka proposisi ``Lulus'' dapat ditulis:
\[
\text{Lulus} \leftrightarrow (p \wedge q).
\]
Tabel kebenaran untuk $p \wedge q$: (B,B)$\to$B, (B,S)$\to$S, (S,B)$\to$S, (S,S)$\to$S. Jadi lulus hanya ketika kedua kondisi terpenuhi.

\begin{contoh}
\textbf{Implementasi kondisi kelulusan dalam Python}

\begin{lstlisting}[style=pythonstyle, caption={Kondisi kelulusan sebagai proposisi majemuk}]
def lulus(ipk, ada_nilai_d):
    """Lulus iff IPK >= 2.0 dan tidak ada nilai D."""
    ipk_cukup = ipk >= 2.0
    tidak_ada_d = not ada_nilai_d
    return ipk_cukup and tidak_ada_d

# Contoh
print(lulus(3.0, False))  # True  -> lulus
print(lulus(2.5, True))   # False -> tidak lulus (ada D)
print(lulus(1.8, False))  # False -> tidak lulus (IPK < 2.0)
\end{lstlisting}
\end{contoh}

Kedua studi kasus di atas menunjukkan bahwa aturan bisnis dan spesifikasi sistem sering dapat dinyatakan sebagai proposisi majemuk dengan operator $\wedge$, $\vee$, $\neg$, $\rightarrow$, dan $\leftrightarrow$, lalu diimplementasikan dengan operator Boolean \texttt{and}, \texttt{or}, \texttt{not} dalam program \cite{python_docs_boolean,rosen2019}.
